\clearpage\begingroup\NoSourceHead\SourceLeaf{330}
\setlength{\fboxsep}{9pt}
\noindent\fbox{\fbox{\begin{minipage}[c][7.5in][c]{\dimexpr\textwidth-4\fboxsep-4\fboxrule\relax}
\begin{center}
{\Large\bfseries CARL FRIEDRICH GAUSS\par}
\vspace{.55in}
{\Huge\bfseries WERKE\par}
\vspace{.55in}
{\large\bfseries ZEHNTEN BANDES ZWEITE ABTEILUNG\par}
\medskip ABHANDLUNG IV
\vspace{.55in}

{\small HERAUSGEGEBEN\par}
\vspace{.3in}
{\small VON DER\par}
\vspace{.3in}
{\large GESELLSCHAFT DER WISSENSCHAFTEN\par}
\vspace{.3in}
{\small ZU\par}
\vspace{.3in}
{\large\bfseries\textls[180]{GÖTTINGEN}\par}
\vspace{.65in}
IN KOMMISSION BEI B. G. TEUBNER IN LEIPZIG\par
\medskip 1923.
\end{center}
\end{minipage}}}
\clearpage\SourceLeaf{332}\NoSourceHead
\HistoricalFoot{X\textsubscript{2} Abh. 4.}{1}
\begin{center}
\vspace*{.5in}
{\LARGE GAUSS ALS GEOMETER\par}
\vspace{1in}
VON
\vspace{.6in}

{\large PAUL STÄCKEL\par}
\vfill
\rule{\linewidth}{.4pt}\par
{\small Abdruck aus Heft 5 der \textit{Materialien für eine wissenschaftliche Biographie von Gauss}\par
gesammelt von F. \textsc{Klein}, M. \textsc{Brendel} und L. \textsc{Schlesinger}.\par
Nachrichten der K. Gesellschaft der Wissenschaften zu Göttingen. Mathem.-physik. Klasse. 1917.\par
Vorgelegt in der Sitzung vom 26. Oktober 1917.\par}
\rule{\linewidth}{.4pt}
\end{center}
\clearpage\SourceLeaf{334}\NoSourceHead\HistoricalFoot{}{1\textsuperscript{*}}
\renewcommand{\thefootnote}{\arabic{footnote})}\setcounter{footnote}{0}
\begin{center}\textbf{1.}\par\medskip\textbf{\textls[80]{Einleitung}\footnotemark[1].}\end{center}
\footnotetext[1]{\textls[100]{Verzeichnis der Abkürzungen.}\par
\noindent W. für C. F. \textsc{Gauss}, Werke I—XI.\par
\noindent T. für das Wissenschaftliche Tagebuch, W. X\textsubscript{1}, S. 488—572.\par
\noindent Br. G.-\textsc{Sch.} für Briefwechsel zwischen \textsc{Gauss} und \textsc{Schumacher} I—VI, Altona 1860—1865.\par
\noindent Br. G.-O. für Briefwechsel zwischen \textsc{Gauss} und \textsc{Olbers}, in W. \textsc{Olbers}, Sein Leben und seine Werke II\textsubscript{1} und II\textsubscript{2}, Berlin 1900 und 1909.\par
\noindent Br. G.-\textsc{Bessel} für Briefwechsel zwischen \textsc{Gauss} und \textsc{Bessel}, Leipzig 1880.\par
\noindent Br. G.-\textsc{Bolyai} für Briefwechsel zwischen \textsc{Gauss} und W. \textsc{Bolyai}, Leipzig 1899.\par
\noindent P. Th. für P. \textsc{Stäckel} und F. \textsc{Engel}, Die Theorie der Parallellinien von \textsc{Euklid} bis auf \textsc{Gauss}, eine Urkundensammlung zur Vorgeschichte der nichteuklidischen Geometrie, Leipzig 1895.\par
\noindent \textsc{Bol.} für W. und J. \textsc{Bolyai}, Geometrische Untersuchungen herausgegeben von P. \textsc{Stäckel}; I. Leben und Schriften der beiden \textsc{Bolyai}, II. Stücke aus den Schriften der beiden \textsc{Bolyai}, Leipzig 1913.\par
\noindent \textsc{Lob.} für N. \textsc{Iw. Lobatschefskij}, zwei geometrische Abhandlungen, aus dem Russischen übersetzt, mit Anmerkungen und mit einer Biographie des Verfassers von F. \textsc{Engel}, Leipzig 1898—99.\par
\noindent \textsc{Sartorius} für W. \textsc{Sartorius v. Waltershausen}, \textsc{Gauss} zum Gedächtniss, Leipzig 1856.\par
\noindent \textsc{Bachmann} für P. \textsc{Bachmann}, Über \textsc{Gauss}' zahlentheoretische Arbeiten, Materialien für eine wissenschaftliche Biographie von \textsc{Gauss}, Heft I, 1911; W. X\textsubscript{2}, Abh. I.\par
\noindent \textsc{Schlesinger} für L. \textsc{Schlesinger}, Über \textsc{Gauss}' Arbeiten zur Funktionentheorie, Materialien usw., Heft III, 1912; W. X\textsubscript{2}, Abh. II.}
\noindent\hspace*{\parindent}\textsc{Gauss} gehört zu den grossen Mathematikern, deren eigentümliche Begabung
schon in der ersten Jugend durch ungewöhnliche Leistungen im
Zahlenrechnen hervortrat. Auch während er das Collegium Carolinum zu
Braunschweig besuchte (1792—1795), hat er viel gerechnet; schon im Jahre
1794 erfand er die Methode der kleinsten Quadrate. Auf umfangreiches
numerisches Beobachtungsmaterial gründen sich auch die 1775 beginnenden
Untersuchungen in der höheren Arithmetik, die 1801 in den \textit{Disquisitiones
arithmeticae} einen ersten Abschluss erhalten. Neben die zahlentheoretischen
Untersuchungen treten in diesen Jahren höchster Schaffenskraft die
\clearpage\SourceLeaf{335}\ArithmeticHead{4}{L}{STÄCKEL, GAUSS ALS GEOMETER.}
\setcounter{footnote}{0}
\noindent Entdeckungen auf dem Gebiete der elliptischen Funktionen, und auch die Algebra
gehört, wie das \textit{Tagebuch}\footnote{Das von \textsc{Gauss} während der Jahre 1799 bis 1815 geführte wissenschaftliche Tagebuch oder Notizenjournal ist abgedruckt W. X\textsubscript{1}, S. 488—572; es wird im Folgenden mit T. angeführt.} zeigt und die Dissertation (1799) bestätigt,
zu den mathematischen Gegenständen, denen sich der junge \textsc{Gauss} zuwendet.
Im Vergleich zur Analysis steht die Geometrie im Hintergrunde; doch lässt
eine Aufzeichnung im \textit{Tagebuch} vom September 1799 (T. Nr. 99) schon die
grosse Frage nach den Gründen der Geometrie anklingen.
\par
Die nun einsetzende astronomische Periode, die sich bis etwa 1816 erstreckt,
bringt nach Aussen hin keine wesentliche Änderung, denn unter den
Veröffentlichungen kommen nur Beiträge zur elementaren Geometrie in Betracht.
Nachlass und Briefwechsel zeigen jedoch, dass die Forschungen über
die Grundlagen der Geometrie nicht geruht haben, und gerade in der Zeit
zwischen 1810 und 1816 ist \textsc{Gauss} zu den grundlegenden Begriffen und Sätzen
aus der Lehre von den krummen Flächen gelangt.
\par
Mit dem Jahre 1816 beginnt die Zeit der Geodäsie. Vorbereitet durch
theoretische Arbeiten über die kürzesten Linien auf dem Sphäroide, betätigt
sich \textsc{Gauss} 1821 bis 1825 bei den Messungen im Felde. Den Weg zu
Grösserem bahnend, verfasst er 1822 die Kopenhagener Preisschrift über die
konforme Abbildung krummer Flächen, und 1828 erscheinen, als reife Frucht
langer Mühen, die \textit{Disquisitiones generales circa superficies curvas}, in denen aus
den Anwendungen heraus ein neuer Zweig der reinen Mathematik selbständiges
Leben gewinnt.
\par
Noch zu einer zweiten Reihe von Untersuchungen hat die geodätische
Tätigkeit den Anstoss gegeben, zu sehr eingehenden Forschungen über die
Grundlagen der Geometrie. Hier ist \textsc{Gauss} nicht dazu gelangt, seine Gedanken
ausführlich niederzuschreiben, und wir sind auf spärliche Notizen und einzelne
Stellen in Briefen angewiesen.
\par
Es folgt die Periode der mathematischen Physik. Als diese etwa 1841
geendet hat, kommt es zu einer Nachblüte der geometrischen Forschung. Es
entstehen die beiden Abhandlungen über Gegenstände der höhern Geodäsie
(1843 und 1846); die Grundlagen der Geometrie werden wieder aufgenommen
und erweiterte Auffassungen gewonnen, geometrische Aufgaben verschiedener
Art werden behandelt, und \textsc{Gauss} kehrt auch zu zwei Gebieten zurück, die
\clearpage\SourceLeaf{336}\ArithmeticHead{5}{R}{EINLEITUNG.}
\setcounter{footnote}{0}
\noindent ihn von jeher angezogen hatten und denen er hohe Bedeutung beimass: zur
Geometria situs und zur geometrischen Versinnlichung der komplexen Grössen.
\par
Wie \textsc{Sartorius}\footnote{\textsc{Sartorius von Waltershausen}, \textit{Gauss zum Gedächtnis}, Leipzig 1856; im Folgenden mit \textsc{Sartorius} angeführt.} berichtet (S. 80), hat \textsc{Gauss} sich dahin geäussert, »in
seiner frühesten Jugend habe ihm die Geometrie wenig Interesse eingeflösst,
welches sich erst später bei ihm in hohem Masse entwickelt habe«. Die
Arithmetik war und blieb ihm die »Königin der Mathematik«, deren Hofstaat
die andern Zweige der Analysis angehörten. Gewiss war ihm das geometrisch-anschauliche
Denken nicht fremd, aber bei seinen geometrischen Untersuchungen
hat er fast überall die analytischen Methoden bevorzugt. »Es ist nicht zu
leugnen«, heisst es in der Besprechung der \textit{Géométrie descriptive} von \textsc{Monge}
(W. IV, S. 359), »dass die Vorzüge der analytischen Behandlung vor der
geometrischen, ihre Kürze, Einfachheit, ihr gleichförmiger Gang, und besonders
ihre Allgemeinheit, sich gewöhnlich um so entschiedener zeigen, je
schwieriger und verwickelter die Untersuchungen sind«. Er war sich jedoch
dessen wohl bewusst, dass »die logischen Hilfsmittel für sich nichts zu leisten
vermögen und nur taube Blüten treiben, wenn nicht die befruchtende, lebendige
Anschauung des Gegenstandes überall waltet« (W. IV, S. 366). Die Pflege
der rein geometrischen Methoden hielt er für »unentbehrlich beim frühern
jugendlichen Studium, um Einseitigkeit zu verhüten, den Sinn für Strenge
und Klarheit zu schärfen und den Einsichten eine Lebendigkeit und Unmittelbarkeit
zu geben, welche durch die analytischen Methoden weit weniger befördert,
mitunter eher gefährdet werden« (W. IV, S. 360), und er wünschte,
»dass auch die rein geometrischen Behandlungen fortwährend kultiviert werden
und dass die Geometrie wenigstens einen Teil der neuen Felder, die die Analyse
erobert, sich aneigne« (W. II, S. 186).
\begin{center}\rule{.18\linewidth}{.4pt}\end{center}
\clearpage\endgroup
